% Generated draft: requires mathematical and source review.
\documentclass[11pt]{article}
\usepackage[margin=25mm]{geometry}
\usepackage{fontspec}
\setmainfont{Latin Modern Roman}
\usepackage{amsmath,amssymb,mathtools,unicode-math}
\setmathfont{Latin Modern Math}
\usepackage{enumitem,needspace,xurl,hyperref}
\setlength{\parindent}{0pt}
\setlength{\parskip}{5pt}
\setlength{\emergencystretch}{3em}
\newcommand{\N}{\mathbb{N}}
\newcommand{\Z}{\mathbb{Z}}
\newcommand{\Q}{\mathbb{Q}}
\newcommand{\R}{\mathbb{R}}
\newcommand{\C}{\mathbb{C}}
\newcommand{\F}{\mathbb{F}}
\newcommand{\A}{\mathbb{A}}
\newcommand{\PP}{\mathbb{P}}
\newcommand{\E}{\mathbb{E}}
\newcommand{\eps}{\varepsilon}
\newcommand{\dd}{\,\mathrm{d}}
\begin{document}
\section*{UnsolvedMath ID: 30003309}
\textbf{Source problem code:} catalogue entry
\subsection{Definitions and mathematical statement}
For a field \(k\), an integer \(n\geq 1\), and positive integers \(a_1,\ldots,a_n\), consider the homomorphism
\[
\varphi_{\mathbf a}:k[x_1,\ldots,x_n]\longrightarrow k[t],\qquad x_i\longmapsto t^{a_i}.
\]
Under the standard positive-exponent interpretation, the associated affine monomial curve is
\[
C_{\mathbf a}=\operatorname{Spec}k[t^{a_1},\ldots,t^{a_n}]\subseteq\mathbb A_k^n,
\]
and its defining ideal is \(I_{\mathbf a}=\ker\varphi_{\mathbf a}\). The curve is called Gorenstein when its coordinate ring is Gorenstein: every localization at a prime ideal is a Gorenstein local ring, meaning a Noetherian local ring of finite injective dimension as a module over itself. Define the minimal number of ideal generators by
\[
\mu(I_{\mathbf a})=\min\bigl\{r\in\mathbb Z_{\geq0}:I_{\mathbf a}=(f_1,\ldots,f_r)\text{ for some }f_i\in k[x_1,\ldots,x_n]\bigr\}.
\]
\textbf{Question, retaining the supplied universal quantifier.} For a fixed field \(k\), does there exist an integer \(B\geq0\), independent of both \(n\) and the curve, such that
\[
\forall n\geq1\ \forall(a_1,\ldots,a_n)\in\mathbb Z_{>0}^{\,n},\qquad
k[t^{a_1},\ldots,t^{a_n}]\text{ Gorenstein}
\ \Longrightarrow\ \mu(I_{\mathbf a})\leq B?
\]
This is a conditional formalization using the usual positive-exponent convention. The supplied record does not specify the ground field, whether the parametrization must be minimal, or whether the bound must also be independent of the field. No minimality restriction is imposed in the displayed formulation. In particular, the phrase ``independently of \(n\)'' has been retained literally; its compatibility with the intended source question requires mathematical review. The supplied open-status assessment is not an independently verified conclusion.
\subsection{Short English statement}
The supplied question asks whether one constant bounds the minimum number of equations generating the defining ideal of every affine Gorenstein monomial curve, uniformly over all ambient dimensions. The ground-field convention, permitted parametrizations, and intended uniformity in dimension require source review.
\subsection{Sources}
\begin{itemize}
\item Supplied local registry record: \texttt{data/unsolvedmath/problems.json}, permanent UnsolvedMath ID \(30003309\), source problem code \texttt{catalogue entry}, title ``Generator Bounds for Gorenstein Monomial Curve Ideals.'' The record assigns proposed year \(2016\) and names no proposer.
\item Supplied report citation: \emph{Asymptotic Phenomena in Local Algebra and Singularity Theory} (2017), Oberwolfach Reports, pp.~3199--3259, DOI \texttt{10.4171/owr/2016/57}. \url{https://doi.org/10.4171/owr/2016/57}. The collection describes these reports as documenting mathematical workshops and their open-problem discussions.
\item Supplied literature source, labeled ``Bresinsky, monomial Gorenstein ideals'': \url{https://doi.org/10.1007/BF01303625}. No fuller bibliographic description is supplied.
\item Supplied literature source, labeled ``OWR 57 (2017)'': \url{https://doi.org/10.4171/owr/2016/57}.
\item Literature assessment dated exactly \texttt{2026-08-21}: the supplied curation labels the problem open and reports bounds or unbounded families in restricted settings while describing the universal boundedness question as open. This is supplied dated curation, derived from the supplied OWR curation, and has not been independently verified; its status and cited claims require independent checking.
\item The supplied statement status is \texttt{corrected\_verified}. Its accompanying verification note states that only safe notation, typography, and TeX-delimiter cleanup was applied, with hypotheses, quantifiers, and subquestions preserved. This label does not establish mathematical correctness or independent verification here.
\end{itemize}
\end{document}
